\documentclass{article}
\usepackage[margin=1in]{geometry}
\usepackage{mathtools, amsfonts, amsthm}

\usepackage{psvectorian} % This is the only specific package you need
\usepackage[svgnames]{xcolor}

\title{HW 0}
\author{Sam Ly}

\begin{document}
\maketitle

\section*{Exercise 1 [2 pts]}

I posted an image of my music player and headphones into the 
Discord introduction channel. Finding the assignment on Gradescope
was straightforward enough.

\pagebreak

\section*{Exercise 2 [8 pts]}

\begin{enumerate}
    \item[i.] {
        I already feel comfortable with \LaTeX.
    }
    \item[iv.] {
        First, when using enumerate tag, for each item, using the 
        square braces lets you define custom "numbers" for the list. Here,
        we use this feature to use Roman numerals and skip some numbers.
        \begin{verbatim}
            \begin{enumerate}
                \item[i.] First.
                \item[iii.] Third!
            \end{enumerate}
        \end{verbatim}

        Then, you can use `.sty' files to add whimsical ornaments to your
        document! This is from the \begin{verbatim}
            \usepackage{psvectorian}
        \end{verbatim}
        package.

        \psvectorian[width=0.25\textwidth]{105}
        \psvectorian[width=0.4\textwidth]{84}
        \psvectorian[width=0.25\textwidth]{106}

        Not sure how well the discord bot will handle this.
    }
\end{enumerate}

\section*{Exercise 3 [4 pts]}

\begin{enumerate}
    \item {
        Watched.
    }
    \item {
        Although I was able to get the ``gist'' of the video, some details left 
        me a little bit confused. The first is what he means at 14:30, when he 
        says that certain groups can be ``factored'' into smaller groups. I get 
        that at some point, groups can be viewed as their own objects, but for 
        now, I struggle to conceptualize the act of "composing" groups themselves.
        It feels like there is some higher level of abstraction that I am not yet 
        familiar with. Another thing I am curious about is what the ``families''
        of groups entail. What properties constitute a family? He explains some 
        basic examples, but I want to see more of what families are.
    }
\end{enumerate}

\section*{Exercise 4 [3 pts]}

Although I have already read this article, I found some new ideas that I will 
try to implement more heavily. First, I want to get better at deciding what is 
important to say. Although I feel as though I \emph{can} get my ideas accross,
sometimes I feels like I am not necessarily doing it in the most effective way.
Especially in a more abstract (literally) class, where abstractions are literally
the point, I want to get better at dicerning the level of abstract I should 
explain to. Secondly, I want to think more deeply about the structure of my 
arguments. I feel as though the arguments in this class may grow longer, increasing
the importance of good structure.


\section*{Choice Exercise 5 [3 pts]}

\begin{enumerate}
    \item[a.] {
        I understood most of what I went through. For example, the process of 
        listing out the composition of rotations was relatively straightforward.
    }
    \item[b.] {
        One thing I didn't quite fully grasp is page 9. Although I can see that 
        rotations can compose together to form a specific rotation that we have 
        seen, I struggle to see the exact formula for the square and triangle. 
    }
    \item[c.] {
        One thing I am curious about is the impact of spacial dimensions when counting
        symmetries. For example, I noticed that rotations about axis that are 
        parallel to the plain of the shape itself in 3D is analogous to flips 
        in 2D. I feel as though this is significant,
        as it may allow us to conceptualize higher dimensional symmetries by 
        building lower dimensional analogs.
    }
\end{enumerate}

\section*{Choice Exercise 7 [5 pts]}

\begin{enumerate}
    \item[2.] {
        \begin{enumerate}
            \item {
                Let

                \[A =
                    \begin{bmatrix}
                    a^1_1 & a^1_2 & \cdots & a^1_n \\
                    a^2_1 & a^2_2 & \cdots & a^2_n \\
                    \vdots & \vdots & \ddots & \vdots \\
                    a^m_1 & a^m_2 & \cdots & a^m_n
                    \end{bmatrix},
                B = 
                \begin{bmatrix}
                    
                    b^1_1 & b^1_2 & \cdots & b^1_n \\
                    b^2_1 & b^2_2 & \cdots & b^2_n \\
                    \vdots & \vdots & \ddots & \vdots \\
                    b^m_1 & b^m_2 & \cdots & b^m_n
                \end{bmatrix}.
                \]

                Then, we have \(C = AB\) where
                \[C = 
                \begin{bmatrix}
                    
                    c^1_1 & c^1_2 & \cdots & c^1_n \\
                    c^2_1 & c^2_2 & \cdots & c^2_n \\
                    \vdots & \vdots & \ddots & \vdots \\
                    c^m_1 & c^m_2 & \cdots & c^m_n
                \end{bmatrix}.
                \]
                such that \(c^i_j = \sum{a^i_k b^k_j}\).
            }
            \item {
                To show that the the operation of matrix multiplication is 
                associative, we must show that \(A(BM) = (AB)M\), where 
                \(A, B,\) and \(M\) are \(n\times n\) matrices. By the definition
                of matrix multiplication, we see that 
                \[(AB)^i_j = \sum_k^n{a^i_kb^k_j}.\]
                Then, we see that
                \[((AB)M)^i_j = \sum_k^n{(AB)^i_k m^k_j}.\]
                By expanding \((AB)^i_j\), we get 
                \[((AB)M)^i_j = \sum_k^n{\sum_l^n{a^i_l b^l_j} m^k_j}.\]

                Then, we perform similar manipulations for \(A(BM)\) to see that
                \[(BM)^i_j = \sum_k^n{b^i_k m^k_j}.\]

                Then, we see that 
                \[(A(BM))^i_j = \sum_l^n{a^i_l (BM)^i_j}.\]

                By expanding \((BM)^i_j\), we see 
                \[(A(BM))^i_j = \sum_l^n{a^i_l \sum_k^n{b^i_k m^k_j}},\]
                
                which can be rewritten as
                \[(A(BM))^i_j = \sum_l^n{\sum_k^n a^i_l b^i_k m^k_j}.\]

                Thus, \((A(BM) = (AB)M\) for any \(n \times n\) matrices \(A, B, C\),
                so matrix multiplication is associative.
            }
        \end{enumerate}
    }
\end{enumerate}

\end{document}
